% Source-aligned English for printed p.122.
% Decisions: TOH-P122-R1-EN-001 through TOH-P122-R1-EN-010.
\SourcePage{122}

\providecommand{\TohokuContentsChapter}[1]{%
  \par\medskip\noindent\textbf{#1}\par}
\providecommand{\TohokuContentsEntry}[2]{%
  \par\noindent\hangindent=2.8em\hangafter=1%
  \makebox[2.8em][l]{#1}#2\par}

\LayerTarget{tohoku:contents}
\begingroup\small
\LayerTarget{tohoku:section:4.2}
\TohokuContentsEntry{}{fundamental}
\LayerTarget{tohoku:section:4.3}
\TohokuContentsEntry{4.3.}{Case of a constant sheaf of rings}
\LayerTarget{tohoku:section:4.4}
\TohokuContentsEntry{4.4.}{Case of sheaves with a group of operators}

\LayerTarget{tohoku:chapter:V}
\TohokuContentsChapter{CHAPTER V\quad COHOMOLOGICAL STUDY OF SPACES WITH
OPERATORS}

\LayerTarget{tohoku:section:5.1}
\TohokuContentsEntry{5.1.}{Generalities on $G$-sheaves}
\LayerTarget{tohoku:section:5.2}
\TohokuContentsEntry{5.2.}{The functors $H^{n}(X;\,G,A)$ and
$\mathbf{H}^{n}(G,A)$ and the fundamental spectral sequences}
\LayerTarget{tohoku:section:5.3}
\TohokuContentsEntry{5.3.}{Case of a discontinuous group of homeomorphisms}
\LayerTarget{tohoku:section:5.4}
\TohokuContentsEntry{5.4.}{Transformation of the first spectral sequence}
\LayerTarget{tohoku:section:5.5}
\TohokuContentsEntry{5.5.}{Calculation of the $H^{n}(X;\,G,A)$ by coverings}
\LayerTarget{tohoku:section:5.6}
\TohokuContentsEntry{5.6.}{The functors
$\mathrm{Ext}_{O,G}^{n}(X;\,A,B)$}
\LayerTarget{tohoku:section:5.7}
\TohokuContentsEntry{5.7.}{Introduction of families $\Phi$.}
\endgroup

\LayerTarget{tohoku:chapter:I}
\LayerContents{Chapter I — Generalities on abelian categories}
\begin{center}
  \textbf{Chapter I\quad Generalities on abelian categories}
\end{center}

\LayerTarget{tohoku:section:1.1}
\noindent\textbf{1.1. Categories.}\quad
\LayerTarget{tohoku:section:1.1:p1}
Recall that a \textit{category} is a nonempty class $C$ of objects, together
with, for $A,B\in C$, a set $\mathrm{Hom}(A,B)$ (called the set of
\textit{morphisms} from $A$ to $B$), and, for three objects $A,B,C\in C$, a
map (called \textit{composition of morphisms})
$(u,v)\longrightarrow vu$ from
$\mathrm{Hom}(A,B)\times\mathrm{Hom}(B,C)$ to $\mathrm{Hom}(A,C)$, these data
satisfying the following two axioms: composition of homomorphisms is
\textit{associative}; for every $A\in C$, there exists in
$\mathrm{Hom}(A,A)$ an element $I_A$ (called the \textit{identity morphism of
$A$}) that is a right and left identity for composition of morphisms. (This
element $I_A$ is then uniquely determined.) Finally, it will be prudent to
suppose that the datum of a morphism $u$ determines its ``source'' and
``target'' objects, in other words, that if $(A,B)$ and $(A',B')$ are two
distinct pairs of objects of $C$, then $\mathrm{Hom}(A,B)$ and
$\mathrm{Hom}(A',B')$ are two disjoint sets.

\LayerTarget{tohoku:section:1.1:p2}
Given a category $C$, define the \textit{dual category} $C^{\circ}$ to be the
category having the same objects as $C$, and in which the set
$\mathrm{Hom}(A,B)^{\circ}$ of morphisms from $A$ to $B$ is identical with
$\mathrm{Hom}(B,A)$, the composite of $u$ and $v$ in $C^{\circ}$ being
defined as the composite of $v$ and $u$ in $C$. Every notion and every
statement concerning an unspecified category has a \textit{dual notion or
statement} (``reversal of arrows''), which will be just as useful in
applications, but whose explication is most often left to the reader.

\LayerTarget{tohoku:section:1.1:p3}
Let a category $C$ and a morphism $u:A\longrightarrow B$ in $C$ be given. For
every $C\in C$, define a map
$v\longrightarrow uv:\mathrm{Hom}(C,A)\longrightarrow\mathrm{Hom}(C,B)$ and
a map
$w\longrightarrow wu:\mathrm{Hom}(B,C)\longrightarrow\mathrm{Hom}(A,C)$.
We say that $u$ is a \textit{monomorphism}, or that $u$ is \textit{injective}
(respectively, that $u$ is an \textit{epimorphism}, or that $u$ is
\textit{surjective}), if the first (respectively, the second) of the two
preceding maps is always injective; $u$ is called \textit{bijective} if it is
both injective and surjective.

\LayerTarget{tohoku:section:1.1:p4}
A \textit{left inverse} (respectively, \textit{right inverse}) of $u$ is a
$v\in\mathrm{Hom}(B,A)$ such that $vu=I_A$ (respectively, $uv=I_B$); $v$ is
called an \textit{inverse of $u$} if it is both a left inverse and a right
inverse (in which case it is uniquely determined). $u$ is called an
\textit{isomorphism} if it admits an inverse. If $u$ admits a left inverse
(respectively, a right inverse), it is injective (respectively, surjective),
so an isomorphism is

\TohokuContinuation{123}
